\section{选读：图神经网络、强化学习与多模态研究}
\label{ch:lower-ml-frontiers}

本节的可操作实验限于三节点传播和两期动态规划，分别体现关系共享与行动反馈。它们不训练大型 GNN 或深度强化学习模型。


这一节提供后续学习入口。基础目标仅为识别图上的邻域聚合、序贯决策的状态与动作，以及多模态数据的不同输入；不要求初学者凭本节掌握 GCN、GraphSAGE、GAT、DQN 或 PPO 的完整训练理论。可在学完树模型与验证章节后跳过本节，继续 Alpha 决策。需要深入时，先选一个问题，再学习对应算法、完整训练实验与专门教材。

\MFQLead{图表示与消息传递}

图 $G=(V,E)$ 由节点、边和可选属性组成。量化研究中的节点可以是公司、基金、行业或账户；边可以来自供应链、共同持股、行业隶属、新闻共现或历史相关。前四类更接近外生关系，相关边则由数据估计，必须按时间滚动并防止使用未来窗口。

% Generated by tools/build_figures.py; edit figures/figure-specs.json instead.
\MFQFigure{tex/figures/generated/lower-ml-gnn.pdf}{GNN 通过邻居消息的聚合与更新传播关系信息；每条边的含义和可得时间决定它能提供什么信息。}{lower-ml-gnn}


消息传递 GNN 的一般形式为
\[
m_v^{(l)}=\operatorname{AGG}^{(l)}
\{\psi^{(l)}(h_v^{(l)},h_u^{(l)},e_{uv}):u\in\mathcal N(v)\},
\]
\[
h_v^{(l+1)}=\phi^{(l)}(h_v^{(l)},m_v^{(l)}).
\]
聚合必须对邻居顺序不敏感，常用 sum、mean 或 attention。$L$ 层后节点能接收 $L$ 跳邻域信息；层数过深会造成 over-smoothing，使节点表示趋同。

\MFQTransition{GCN、GraphSAGE 与 GAT}

GCN 常写成
\[
H^{(l+1)}=\sigma\!\left(
\tilde D^{-1/2}\tilde A\tilde D^{-1/2}H^{(l)}W^{(l)}
\right),
\]
其中 $\tilde A=A+I$。归一化防止高阶节点简单主导聚合。GraphSAGE 对邻居采样并学习聚合器，适合大图和新节点；GAT 用注意力学习邻居权重，但权重仍不是因果影响。

\MFQTransition{三节点手算}

考虑链 $1-2-3$，每个节点加自环，标量特征为 $(1,2,4)$。若暂用简单均值聚合，第一层得到
\[
h_1'=1.5,\qquad h_2'=\frac{1+2+4}{3}=\frac73,
\qquad h_3'=3.
\]
节点 2 同时吸收两端信息。若边 $2-3$ 在预测日后才出现，把当前完整图用于历史训练就会泄漏。图快照必须像价格数据一样带生效时间。

图任务包括节点回归（预测个股收益/风险）、边预测（关系形成）、图分类（组合或市场状态）和异构图学习。供应商、客户、行业和高管是不同边类型，强行压成无类型邻接会丢失语义。关系方向、权重、发布日期和删除时间都要进入数据协议。

异构图可为每种关系 $r$ 使用独立变换：
\[
h_v^{(l+1)}=\phi\!\left(
W_0h_v^{(l)}+\sum_r\sum_{u\in\mathcal N_r(v)}
\frac{1}{c_{v,r}}W_rh_u^{(l)}
\right).
\]
它能区分“客户--供应商”和“同属行业”，但参数量随关系类型增加。样本稀疏时，关系专属矩阵可能只记住少数大公司；可以共享低秩基、合并经济含义相近的关系，或把关系属性作为消息输入。无论采用哪种方法，都要报告每类边的覆盖率、更新时间和孤立节点比例。

动态图需要区分事件发生时间、数据库入库时间和研究者首次可见时间。若边在 $t$ 时刻出现，预测 $t$ 之前节点时不能让该边参与消息传递。常见实现是为每个决策日构造快照 $G_t$，或让消息携带时间编码并仅聚合 $t_{uv}\le t$ 的边。仅按最终图切分节点无法阻止关系泄漏，因为训练节点会通过未来边把测试期信息传播回来。

链路预测常用负采样，但“没有记录的边”不一定是真负例，可能只是数据源未覆盖。随机负采样还会生成经济上不可能的公司组合，使任务虚假简单。更可靠的评估按行业、规模和时间匹配候选负边，并分别报告新边发现、旧边持续和删除边识别。AUC 很高但只能识别数据覆盖差异的模型，没有形成可交易关系信号。

GNN 还面对 over-squashing：远处大量信息被压入固定维度向量，梯度和信号都可能衰减。增加层数同时加剧 over-smoothing 与过拟合。残差、跳连、图重连和子图采样可以缓解数值问题，却不能替代“究竟需要几跳经济关系”的研究假设。每增加一跳，都应做删除该跳关系的消融并检查样本外稳定性。

\MFQTransition{强化学习：预测问题变成控制问题}

马尔可夫决策过程由 $(\mathcal S,\mathcal A,P,r,\gamma)$ 构成。策略 $\pi(a\mid s)$ 的价值为
\[
V^\pi(s)=\mathbb E_\pi\left[\sum_{t=0}^{\infty}
\gamma^t r_{t+1}\mid S_0=s\right],
\]
动作价值满足 Bellman 方程
\[
Q^\pi(s,a)=\mathbb E\left[r_{t+1}+\gamma
\mathbb E_{a'\sim\pi}Q^\pi(S_{t+1},a')\mid s,a\right].
\]
在交易中，状态可能包含价格、库存和剩余时间，动作是订单或目标仓位，奖励必须扣除成本、风险和约束罚项。若奖励只用账面收益，代理会通过极端杠杆或不成交订单钻空子。

% Generated by tools/build_figures.py; edit figures/figure-specs.json instead.
\MFQFigure{tex/figures/generated/lower-ml-drl.pdf}{策略选择动作，动作影响状态与奖励，奖励再用于改进策略；离线历史数据没有直接给出未执行动作的后果。}{lower-ml-drl}


\MFQTransition{两步执行算例}

需在两期买入 2 单位，第一期最多买 1 单位，第二期最多买 2 单位。第一期立即买 1 单位成本为 1，等待到第二期买剩余单位的期望成本为 3；若第一期全部等待，第二期买 2 单位期望成本为 6。无折扣时，第一期买 1 的总期望成本为 $1+3=4$，优于等待的 6。这个小型动态规划是复杂 RL 环境的 解析参照：训练出的策略若连它都不能恢复，不能进入市场仿真。

\MFQTransition{DQN、策略梯度与 Actor--Critic}

DQN 用神经网络近似 $Q(s,a)$，目标为
\[
y=r+\gamma\max_{a'}Q_{\theta^-}(s',a'),
\qquad L=(Q_\theta(s,a)-y)^2.
\]
经验回放打破相邻样本相关，目标网络降低自举目标的快速漂移。连续动作更适合策略梯度或 actor--critic。策略梯度定理给出
\[
\nabla_\theta J(\theta)=
\mathbb E\big[\nabla_\theta\log\pi_\theta(a_t\mid s_t)
A^\pi(s_t,a_t)\big].
\]
PPO 用裁剪概率比限制单次更新，但不能消除环境错设和样本外制度变化。

金融 RL 的难点不是调用算法，而是环境反事实不可观察：历史只记录实际采取的一条轨迹，换一个订单会改变成交和市场冲击。纯历史回放无法准确评价未执行动作。离线 RL 还面对行为策略覆盖不足；策略选择了数据中从未出现的动作时，Q 值可能任意外推。

若历史由行为策略 $\mu$ 产生，目标策略 $\pi$ 的轨迹重要性权重为
\[
w_{0:T}=\prod_{t=0}^{T}
\frac{\pi(a_t\mid s_t)}{\mu(a_t\mid s_t)}.
\]
普通重要性采样用 $w_{0:T}$ 加权回报，在策略差异或期限较长时方差会爆炸；逐步重要性采样、加权归一化和 doubly robust 估计可降低方差，但都要求行为概率可估且具有支持覆盖。若 $\mu(a\mid s)=0$ 而 $\pi(a\mid s)>0$，仅凭该离线数据无法识别目标动作价值。

保守离线 RL 在拟合 Bellman 目标之外，惩罚数据外动作的高 Q 值；行为克隆约束则限制策略偏离历史行为。这能减少外推风险，也可能保留历史策略的缺陷。改变保守程度，并与简单规则策略比较，可以观察估计价值、行为偏离和数据支持范围之间的取舍。

交易奖励通常是多目标的。一个执行奖励可写成
\[
r_t=-x_t(P_t-S_0)-\lambda_q q_t^2-\lambda_u u_t^2
-\lambda_f\mathbf 1\{t=T,\ q_t\ne0\},
\]
其中 $q_t$ 是剩余库存，$u_t$ 是本期动作。各项单位必须统一为现金或基点。改变 $\lambda_q$ 不只是“调参”，而是在改变风险偏好；应绘制成本、完成率、尾部损失和库存暴露的 Pareto 前沿，而不是只选总奖励最大的种子。

部分可观测市场更接近 POMDP：代理只能从订单簿切片、成交和延迟推断隐含流动性状态。堆叠最近若干帧、使用循环状态或贝叶斯 belief 都是近似。若模拟器把真实隐状态直接提供给代理，得到的策略使用了额外信息。它可以作为理想信息情形的比较对象，但与仅使用可观测量的策略回答不同的问题。

\begin{diagnostic}
固定历史价格路径的回放，可用于研究价格接受者的决策。执行和做市问题还涉及动作对成交、库存及市场冲击的影响；忽略这些反馈，会改变待求解的控制问题。因此需要比较不同成交与冲击假设下的策略表现。
\end{diagnostic}

\MFQTransition{自监督学习与表示预训练}

未标注金融数据远多于可靠收益标签。自监督任务可以遮蔽重建、下一事件预测、对比学习或跨模态对齐。对比损失将同一对象的增强视图拉近、不同对象推远：
\[
L_i=-\log\frac{\exp(\operatorname{sim}(z_i,z_i^+)/\tau)}
{\sum_j\exp(\operatorname{sim}(z_i,z_j)/\tau)}.
\]
关键在于增强是否保持经济语义。随意打乱时间、缩放价格或替换行业可能改变标签含义。预训练窗口也必须早于下游评价期。

Masked autoencoder 可在收益/盘口序列上重建遮蔽片段；预测任务过容易时，模型只学局部插值。应使用线性 probe、事先确定编码器和完整微调三档比较，判断价值来自表示还是下游容量。

不同预训练目标强调不同关系：下一事件预测利用局部动态，遮蔽重建利用输入窗口内的上下文，对比学习利用增强前后应保持的特征。双向编码器能够联合读取窗口两端的信息；用于在线预测时，整个输入窗口都应已在决策前可得。比较从头训练、冻结编码器和端到端微调时，还需使用一致的语料截止与下游划分，才能区分预训练信息和参数更新的作用。

对比学习还可能发生表示塌缩：所有样本映射到近似同一点。负样本、停止梯度、方差正则或教师--学生更新可避免平凡解，但会引入新的批大小和动量选择。研究报告至少给出表示维度的方差、协方差谱、最近邻语义和下游稳定性；只报告较低预训练损失无法证明表示有用。

\MFQTransition{多模态：数值、文本、图和图像如何对齐}

公司研究同时包含财务表、价格序列、公告文本、供应链图和电话会音频。早期融合直接拼接表示，晚期融合分别预测后再组合，cross-attention 让一个模态查询另一个模态。最难的是实体和时间对齐：公告属于哪家公司、何时首次可见、财务数据何时发布、文本修订何时发生。

缺失模态不能简单补零。训练时可做 modality dropout，推理时报告可用模态集合。若某模态只在大型公司可得，模型可能把“数据可得性”当成规模信号。

跨模态对齐常以配对损失学习共同空间。例如文本表示 $z_i^T$ 与数值表示 $z_i^N$ 的对称对比损失同时让正确配对相互接近：
\[
L=\frac12\bigl[L_{T\to N}+L_{N\to T}\bigr].
\]
配对键必须是“实体--事件--首次可见时点”，而不是只按公司和自然日连接。一天内多份公告、盘前与盘后新闻、跨时区财报都可能造成一对多关系。错误对齐会让大模型学到时间邻近或公司规模，而不是事件语义。

融合模型的消融应包含单模态基线、随机打乱某模态、删除模态和时间错位负例。若加入文本后只在文本覆盖最密集的资产上改善，应进一步区分信息增量和股票池变化。语音转写和文本推理还需要时间。如果输出晚于决策时点，这些特征就无法用于原先定义的预测时刻。

\MFQTransition{生成模型、扩散模型和合成数据}

扩散模型通过逐步加噪和反向去噪学习数据分布，可生成情景或高维序列。GAN 和 VAE 也可用于市场模拟。合成数据的首要评价不是视觉逼真，而是是否保留边际分布、依赖、尾部、条件响应和策略排序，同时不复制训练样本。

生成模型不能凭空提供未观察制度下的因果反事实。用同一生成器训练和评价策略会产生“模拟器过拟合”。至少要跨生成机制评价、保留真实留出期，并明确合成结果只证明在假设环境中的行为。

扩散模型的前向过程逐步加入高斯噪声，反向模型学习噪声或 score：
\[
q(x_t\mid x_{t-1})=
N\!\left(\sqrt{1-\beta_t}x_{t-1},\beta_tI\right),
\qquad
L_{\rm simple}=\mathbb E\|\varepsilon-
\varepsilon_\theta(x_t,t,c)\|_2^2.
\]
条件 $c$ 可以是波动状态、宏观情景或资产属性。生成器若未显式条件于这些变量，抽到的极端路径可能只是在边际分布上罕见，却不对应给定压力情景。反之，条件标签来自未来全样本分类也会泄漏制度信息。

合成市场的验证至少分三层：第一层比较边际、尾部、自相关和横截面相关；第二层比较状态持续、冲击响应与成交机制；第三层比较多种事先确定策略的相对排序及失败模式。没有任何单层指标能证明“足够真实”。策略只在一个生成器上占优时，应把结论写成模拟器条件结论，而不是市场结论。

\MFQTransition{怎样继续学习}
图模型需要先掌握矩阵乘法与监督学习；强化学习需要条件期望、Markov 链与动态规划；生成模型需要概率变换与梯度优化。读者可只选择其中一个问题，先做完下面的小型计算，再进入专门教材与完整训练实验。


\subsection{Bellman 递推与图上的重复平均}
前述两期买入例中，终端成本为 $V_2(R)=3R$，第一期可选 $a\in\{0,1\}$，每单位成本为 1。因此 $V_1(2)=\min\{V_2(2),1+V_2(1)\}=\min\{6,4\}=4$。若允许第一期买 2，答案将变成 2；动作集合是问题的一部分。

三节点均值传播写成 $h'=Ah$，其中
\[
 A=\begin{pmatrix}1/2&1/2&0\\1/3&1/3&1/3\\0&1/2&1/2\end{pmatrix}.
\]
行和为 1 保证每次是邻域凸组合。反复应用会减小节点间差异；在这个连通且有自环的图上，极限由稳态权重 $(2,3,2)/7$ 决定，对初值 $(1,2,4)$ 得共同极限 $16/7$，不是普通均值 $7/3$。归一化方式因此改变传播的含义。notebook 画出每层三个节点的值，并枚举两期的全部合法动作。
\section{现代人工智能方法综合练习}
\begin{enumerate}[leftmargin=*]
  \item \textbf{口述概念：}tensor 的 shape、dtype、device 和梯度状态分别约束什么？
  \item \textbf{口述概念：}因果卷积、RNN、self-attention 和 Transformer 如何表示历史？
  \item \textbf{口述概念：}GNN 和 DRL 分别把普通监督学习问题增加了什么数学对象？
  \item \textbf{笔试推导：}写出两层 MLP 的前向传播与平方损失梯度链，并解释残差连接为何改善梯度路径。
  \item \textbf{笔试推导：}计算核 $(1,-1)$ 对序列 $(1,3,2,5)$ 的因果卷积；再说明普通对称卷积为何可能泄漏。
  \item \textbf{笔试推导：}解释没有位置编码的 self-attention 对时间置换为何不敏感，并手算两 token softmax 权重。
  \item \textbf{笔试推导：}写出 GCN 一层传播和 Bellman 方程，分别指出图快照泄漏与环境错设发生在哪里。
  \item \textbf{数值编程：}把序列反转，比较均值、因果卷积、RNN、Attention 与 Transformer 输出，并验证 padding mask 与 causal mask 的差别。

  \item \textbf{数值编程：}用两层标量网络比较学习率 $0.1$ 与 $0.01$ 的一次更新，解释损失是否下降。
  \item \textbf{数值编程：}构造三节点供应链图，比较静态全图与逐期图快照的预测差异；不得使用预测日后新增的边。
  \item \textbf{研究判断：}审计一个把均值池化随机特征称作“深度序列模型”的研究材料。
  \item \textbf{研究判断：}比较全量微调、只训任务头和 LoRA/adapter 的证据边界，并列出 LLM 历史回测的参数记忆污染。
  \item \textbf{研究判断：}设计一个执行 RL 的最小解析或独立计算参照，并说明为什么固定历史价格回放不足以评价未执行动作。
  \item \textbf{开放项目：}从序列或文本预测、动态图预测、执行决策中选一个问题，比较简单基线与所学模型，并分析成本及一个模型假设变化的影响。
\end{enumerate}

\subsection*{两层网络的全部梯度：基础练习}
上述旧参数下若 $x=0$，哪些梯度为零？

\subsection*{记忆衰减与注意力的数值含义：基础练习}
$h_t=0.5h_{t-1}+x_t$，输入 $(1,0,0)$，求三个隐状态。

\subsection*{Bellman 递推与图上的重复平均：基础练习}
若三个节点初始特征都为 5，一次均值传播后是多少？

